\documentclass[11pt]{article}
\usepackage[T1]{fontenc}
\usepackage{textcomp}
\usepackage[margin=0.75in]{geometry}
\usepackage{amsmath,amssymb,amsthm}
\usepackage{array,booktabs}
\usepackage{hyperref}
\setlength{\parindent}{0pt}
\newtheorem{theorem}{Theorem}
\theoremstyle{definition}
\newtheorem{definition}{Definition}
\title{Flow Proof Artifact: prop-37-power-of-point-theorem}
\author{Generated by \texttt{flow doc proof}}
\date{Flow Proof Book}
\begin{document}
\maketitle
If a point is taken outside a circle and from it two straight lines fall on the circle, one a tangent and the other a secant, the rectangle on the whole secant and the part outside equals the square on the tangent.\\[0.5em]
\textbf{Source.} euclid — Elements, Book III, Proposition 37\\[0.5em]
\bigskip
\noindent{\large \textbf{Derived fact 1.} \textit{If a point is taken outside a circle and from it two straight lines fall on the circle, one a tangent and the other a secant, the rectangle on the whole secant and the part outside equals the square on the tangent} \hfill the Euclidean plane \cdot Euclid Book III \cdot Proposition 37: the rectangle on a secant and its exterior part equals the square on the tangent}\par
\smallskip\noindent\textit{Coordinate: the Euclidean plane \cdot Euclid Book III \cdot Proposition 37: the rectangle on a secant and its exterior part equals the square on the tangent}\par
\smallskip\noindent\textit{Source: euclid — Elements, Book III, Proposition 37}\par
\smallskip\noindent\textit{Built on: proposition 8: the rectangle on a secant and its exterior segment equals the square on the tangent, for Euclid Book III on the Euclidean plane, proposition 9: if the rectangle equals the square then the line is a tangent, for Euclid Book III on the Euclidean plane}\par
\medskip
\noindent\textbf{Goal.} If a point is taken outside a circle and from it two straight lines fall on the circle, one a tangent and the other a secant, the rectangle on the whole secant and the part outside equals the square on the tangent\par
\noindent$\displaystyle \text{power of point rect equals sq}$\par
\medskip
\noindent
\begin{tabular}{@{} >{\raggedright\arraybackslash}p{0.06\textwidth} >{\raggedright\arraybackslash}p{0.47\textwidth} >{\raggedleft\arraybackslash}p{0.41\textwidth} @{}}
\textbf{\#} & \textbf{Proof} & \textbf{Mathematics} \\
\midrule
\textbf{1.} & We prove that proposition 37: the rectangle on a secant and its exterior part equals the square on the tangent for Euclid Book III the Euclidean plane. & \\[0.45em]
\textbf{2.} & Let P = exterior\_point. & \\[0.45em]
\textbf{3.} & Let secant = line\_PABC. & \\[0.45em]
\textbf{4.} & Let tangent = line\_PT. & \\[0.45em]
\textbf{5.} & From step 2, step 3, and step 4, we can deduce that rect(secant, exterior) equals sq(tangent). & $\displaystyle rect(secant, exterior) = sq(tangent)$ \\[0.45em]
\textbf{6.} & We invoke the derived fact governing Euclid Book III the Euclidean plane: proposition 8: the rectangle on a secant and its exterior segment equals the square on the tangent, for Euclid Book III on the Euclidean plane. & \\[0.45em]
\textbf{7.} & We invoke the derived fact governing Euclid Book III the Euclidean plane: proposition 9: if the rectangle equals the square then the line is a tangent, for Euclid Book III on the Euclidean plane. & \\[0.45em]
\textbf{8.} & From step 6 and step 7, this implies power of point rect equals sq. Hence proven. & $\displaystyle \text{power of point rect equals sq}$ \\[0.45em]
\bottomrule
\end{tabular}
\label{thm:Geometry-Euclid Book III-proposition_37_the_rectangle_on_a_secant_and_its_exterior_part_equals_the_square_on_the_tangent}
\par\medskip\hrule
\end{document}
